\documentclass[aspectratio=169,11pt]{beamer}

\usetheme{default}
\usecolortheme{default}
\setbeamertemplate{navigation symbols}{}

\definecolor{deepblue}{HTML}{1B3A5C}
\definecolor{lightgray}{HTML}{F5F5F5}
\setbeamercolor{frametitle}{fg=deepblue}
\setbeamercolor{title}{fg=deepblue}
\setbeamercolor{structure}{fg=deepblue}
\setbeamercolor{block title}{fg=white,bg=deepblue}
\setbeamercolor{block body}{bg=lightgray}

\usefonttheme{professionalfonts}
\setbeamerfont{frametitle}{size=\large,series=\bfseries}
\setbeamerfont{title}{size=\Large,series=\bfseries}

\setbeamertemplate{headline}{%
  \leavevmode\hbox{\color{deepblue}\rule{\paperwidth}{1.5pt}}%
}
\setbeamertemplate{footline}{%
  \leavevmode\hbox to \paperwidth{%
    \color{deepblue}\rule{\paperwidth}{0.5pt}}%
  \vskip-1pt%
  \hbox to \paperwidth{\hfill\tiny\insertframenumber/\inserttotalframenumber\hspace{4pt}\vspace{2pt}}%
}

\usepackage{graphicx,booktabs,amsmath,hyperref}

\title{Backward Test Particle Tracing in MAGE\\[4pt]
{\normalsize Pilot: St.\ Patrick's Day 2013 Storm}}
\author{Y.\ Zhu \and S.\ Bao \and S.\ Sadeghzadeh \and F.\ Toffoletto}
\institute{Rice University}
\date{}

\begin{document}

\begin{frame}
\titlepage
\end{frame}

% ============ Motivation ============
\begin{frame}{Motivation and Method}

\begin{columns}[T]
\begin{column}{0.52\textwidth}
\textbf{Question}: What fraction of RC H$^+$ comes from bubbles?

\vspace{4pt}
\footnotesize
\begin{tabular}{@{}p{2.0cm}p{4.2cm}@{}}
\toprule
\textbf{Study} & \textbf{Finding} \\
\midrule
Sciola 2023 & $\geq$50\% RC energy from bubbles \\
Yang 2015 & $\sim$60\% pressure from bubbles \\
\bottomrule
\end{tabular}

\vspace{6pt}
\normalsize
\textbf{Two approaches compared}:
\begin{enumerate}
\item CHIMP full-orbit (FO) backward trace in GAMERA fields ($\mathbf{E}=-\mathbf{V}\times\mathbf{B}$, Boris pusher)
\item Sina's RCM backtracker ($V_\text{eff} = \Phi + \lambda V_M$, bounce-averaged drift)
\end{enumerate}
\end{column}
\begin{column}{0.45\textwidth}
\footnotesize
\textbf{Pilot setup}:
\begin{itemize}
\item GAMERA-RCM (sp13\_075), 17 March 2013
\item 1{,}000 H$^+$, 10--200\,keV
\item Uniform in $r_\text{eq}=4.2$--6.6\,R$_E$
\item Backward 4\,h: 10:00$\to$06:00\,UT
\item FO: fixed $\Delta t=0.01$\,s, $x_\text{eq}/y_\text{eq}$ from field line trace
\item Sina: $V_\text{eff}$ drift on $(I,J)$ grid, RK4 with 10 substeps/min
\end{itemize}
\vspace{2pt}
\textbf{Bubble criteria}:\\
$\Delta B_z > 15$\,nT on GAMERA eq.\ plane,\\
midnight sector (MLT 21--03),\\
$x<0$, $r>6$\,R$_E$
\end{column}
\end{columns}
\end{frame}

% ============ Sina method ============
\begin{frame}{Sina's RCM Backtracker}

\begin{columns}[T]
\begin{column}{0.50\textwidth}
\textbf{RCM drift formulation}:
\begin{itemize}
\item Effective potential: $V_\text{eff} = V + \lambda \cdot V_M$
\item $V$: electrostatic potential from REMIX (Vasyliunas eq.)
\item $V_M = (\text{bVol})^{-2/3}$: magnetic drift potential (from flux tube volume)
\item $\lambda$: adiabatic invariant, fixed per particle
\item Energy: $K\,[\text{keV}] = |\lambda| \cdot V_M / 1000$
\end{itemize}

\vspace{4pt}
\textbf{Drift velocity} on ionospheric $(I,J)$ grid:
\[
\frac{dI}{dt} = \frac{1}{\text{fac}}\frac{\partial V_\text{eff}}{\partial J}, \qquad
\frac{dJ}{dt} = -\frac{1}{\text{fac}}\frac{\partial V_\text{eff}}{\partial I}
\]
where $\text{fac} = r_i^2 |B_r| \sin\theta \, \frac{d\phi}{dJ}\frac{d\theta}{dI}$ (Jacobian)
\end{column}
\begin{column}{0.48\textwidth}
\textbf{Inductive E field}:

The electrostatic potential $V$ alone does not contain $-\partial\mathbf{A}/\partial t$. However, the mapping $(I,J) \to (x_\text{eq}, y_\text{eq})$ via $X_\text{MIN}(I,J,t)$ changes each timestep because $\mathbf{B}$ changes. This evolution implicitly captures the inductive E field effect.

\vspace{6pt}
\textbf{Validation vs FO}:

\vspace{2pt}
\footnotesize
\begin{tabular}{@{}lrr@{}}
\toprule
 & $\cos\theta$ & err (R$_E$) \\
\midrule
5\,min & 0.98 & 0.6 \\
10\,min & 0.96 & 1.2 \\
20\,min & 0.88 & 3.3 \\
\bottomrule
\end{tabular}
\normalsize

\vspace{4pt}
Speed ratio $\approx$\,0.83 (Sina systematically slower due to bounce-averaging)
\end{column}
\end{columns}
\end{frame}

% ============ Bubble results ============
\begin{frame}{Results: Bubble Contribution to Ring Current}

\begin{columns}[T]
\begin{column}{0.50\textwidth}
\footnotesize
\textbf{Origin classification} (ever-in-bubble during 4\,h trace):\\
$\Delta B_z > 15$\,nT, midnight MLT 21--03, $x<0$, $r>6$\,R$_E$

\vspace{4pt}
\begin{tabular}{@{}lrr@{}}
\toprule
 & \textbf{FO} & \textbf{Sina} \\
\midrule
Bubble & 30\% & 16\% \\
Pre-existing & 15\% & 70\% \\
Convection & 56\% & 15\% \\
\bottomrule
\end{tabular}

\vspace{6pt}
\normalsize
\textbf{By energy}:

\vspace{2pt}
\footnotesize
\begin{tabular}{@{}lrr@{}}
\toprule
Energy & FO bubble & Sina bubble \\
\midrule
10--30\,keV & 52\% & 26\% \\
30--60\,keV & 38\% & 27\% \\
60--120\,keV & 23\% & 6\% \\
120--200\,keV & 6\% & 6\% \\
\bottomrule
\end{tabular}
\end{column}
\begin{column}{0.48\textwidth}
\footnotesize
\textbf{Key findings}:
\begin{itemize}
\item FO sees \textbf{more bubbles} (30\% vs 16\%) because full-orbit gc drift carries particles through bubble channels faster
\item Sina has \textbf{more pre-existing} (70\% vs 15\%) because bounce-averaged drift is slower $\to$ particles don't escape RC in 4\,h
\item Low-energy particles ($<$30\,keV) have highest bubble fraction in both methods --- ExB convection efficiently transports them through bubbles
\item High-energy ($>$120\,keV) bubble fraction converges ($\sim$6\%) --- these particles drift azimuthally too fast to be captured by bubbles
\end{itemize}
\end{column}
\end{columns}
\end{frame}

% ============ Video ============
\begin{frame}{Video: FO vs Sina with $\Delta B_z$ and RCM Boundary}

\centering
\vspace{10pt}
\textbf{4-hour backward trace animation}\\[4pt]
Left: FO (equatorial projection) \quad Right: Sina $V_\text{eff}$\\[2pt]
Background: GAMERA $\Delta B_z$ \quad Green contour: $\Delta B_z = 15$\,nT (bubble boundary)\\
Magenta line: time-varying RCM active domain boundary\\[2pt]
Red = bubble origin \quad Blue = non-bubble

\vspace{10pt}
\footnotesize
\textbf{Download video}:\\[2pt]
\url{https://drive.google.com/drive/folders/backtrac}\\[2pt]
\texttt{fo\_vs\_veff\_bubble15nt\_midnight.mp4}

\vspace{10pt}
\normalsize
\textit{(Play video during presentation)}

\end{frame}

\end{document}
